%!TEX root = ../hutbthesis_main.tex

\chapter{总结与展望}

\section{总结}

本文针对人车模拟器脚本开发效率低、易出错、调试麻烦、环境部署难等问题，完成人车模拟器快速开发验证插件系统的研究、设计、实现和测试。本文以VS
Code扩展框架为基础，创建出包含编码、校验、调试、部署等全过程的开发辅助工具，主要成果如下所示。

第一，完成了插件系统整体的设计。使用分层架构和模块化设计，确定表现层、业务层、数据层之间的职责界限，提出模块化、可扩展、易用、实时、跨平台这五个设计原则，从而使系统结构清楚，容易维护和迭代。

第二，实现八个主要功能模块。分别为代码智能补全、实时错误检测、多场景联动调试、一体化环境打包、状态栏交互，用抽象语法树解析、Levenshtein编辑距离、状态机参数计数、调试适配器协议、绿色便携打包等关键技术来解决专用API没有提示、错误不能提前发现、调试配置重复、环境迁移复杂等问题。

第三，对整个流程进行了系统的测试。经过测试，插件的功能齐全、运行正常，可以明显加快脚本的编写速度、减少API调用出错率、缩短调试和部署的时间，完全可以满足教学、科研、团队合作等真实的使用场景，具有很高的实用性以及推广的价值。

第四，形成了一套可以复用的领域专用IDE插件开发方案。本文提出的这一个架构、接口和思路，可以给其他垂直领域SDK辅助工具的开发提供一定的借鉴。

\section{展望}

尽管插件已经实现了主要的功能，并且达到了预期的效果，但是由于开发周期以及测试环境的限制，系统还存在着可以改进的空间，可以进一步改善的地方有如下几个方面。

第一点就是提高静态检测的准确性。在已经进行的参数数量校验的基础上，对参数类型、返回值类型、取值范围、空值判断等更严格的要求进行语法和语义的校验，提高错误覆盖率以及准确率。

第二，支持自定义与私有 API
扩展。支持用户导入自定义的接口配置文件，使团队内可以使用不同的函数提示、项目级API文档联动等，提高插件的灵活性以及适用性。

第三，改进远程调试验收以及网络兼容能力。改善弱网、跨设备、容器化环境下的调试连接稳定度，增添日志即时上传，远程断点，多机协作调试这些功能。

第四点就是轻量化以及增量打包的优化。使用差量打包、压缩优化、组件裁剪等方式来减小环境包的体积，加快环境包分发、传输和解压的速度，适合大批量部署。

第五，引入 AI
辅助开发能力。通过代码大模型的自动补全、函数生成、注释生成、代码重构等功能，把``辅助提示''提升为``智能开发''。

将来插件可以和仿真平台联动起来，进行脚本在线校验、一键运行、数据可视化、任务自动化等功能扩展，为人车模拟器、智能驾驶仿真开发提供更高效、更智能、更统一的工具支持。

\textbf{参考文献：}

\begin{enumerate}
\def\labelenumi{\arabic{enumi}.}
\item
  陈夏非,焦翊洋,姜雪,等.基于自适应模型预测控制的驾驶模拟器运动提示算法设计与验证{[}J/OL{]}.机械工程学报,1-10{[}2026-04-28{]}.https://link.cnki.net/urlid/11.2187.th.20260225.1756.008.
\item
  孔令智,田毕江,熊昌安,等.汽车驾驶模拟器的应用研究综述{[}J{]}.公路与汽运,2025,41(04):1-9.DOI:10.20035/j.issn.1671-2668.2025.04.001.
\item
  王俊达,卿杜政.柔性可扩展装备体系对抗仿真建模框架研究{[}J{]}.现代防御技术,2020,48(04):122-131.
\item
  Liuzzo M S ,Carmignani N ,Carver L , et al.Commissioning simulations
  tools based on python Accelerator Toolbox{[}J{]}.Journal of Physics:
  Conference Series,2024,2687(3):DOI:10.1088/1742-6596/2687/3/032001.
\item
  杨博,张能,李善平,等.智能代码补全研究综述{[}J{]}.软件学报,2020,31(05):1435-1453.DOI:10.13328/j.cnki.jos.005966.
\item
  许鹏宇,况博裕,苏铓,等.基于大语言模型的自动代码修复综述{[}J{]}.计算机研究与发展,2025,62(08):2040-2057.
\item
  郭家顺.基于LLM的智能代码补全系统设计与性能评估{[}J{]}.软件,2025,46(10):35-37.
\item
  李溪子.基于机器学习算法的智能代码补全工具开发与实现{[}J{]}.软件,2025,46(02):98-100.
\item
  邹佰翰,汪莹,彭鑫,等.重新审视代码补全中的检索增强策略{[}J{]}.软件学报,2025,36(06):2747-2773.DOI:10.13328/j.cnki.jos.007226.
\item
  陈伟,何成万,余秋惠,等.基于CNN和自注意力神经网络的代码补全方法{[}J{]}.计算机工程与设计,2025,46(10):2919-2926.DOI:10.16208/j.issn1000-7024.2025.10.024.
\item
  杨泽洲.基于语言模型的代码构建辅助技术研究{[}D{]}.哈尔滨工业大学,2025.DOI:10.27061/d.cnki.ghgdu.2025.002727.
\item
  许柏炎,蔡瑞初,梁智豪.一种用于代码注释自动生成的语法辅助复制机制{[}J{]}.计算机工程,2021,47(04):92-99.DOI:10.19678/j.issn.1000-3428.0057290.
\item
  王莹莹.基于预训练模型的目标描述代码自动生成研究{[}D{]}.河北师范大学,2024.DOI:10.27110/d.cnki.ghsfu.2024.000468.
\item
  Hostnik M ,Šikonja R M .Retrieval-augmented code completion for local
  projects using large language models{[}J{]}.Expert Systems With
  Applications,2025,292128596-128596.DOI:10.1016/J.ESWA.2025.128596.
\item
  Husein A R ,Aburajouh H ,Catal C .Large language models for code
  completion: A systematic literature review{[}J{]}.Computer Standards
  \& Interfaces,2025,92103917-103917.DOI:10.1016/J.CSI.2024.103917.
\item
  Wannita T ,Chakkrit T ,Yuan-Fang L .Syntax-aware on-the-fly code
  completion{[}J{]}.Information and Software
  Technology,2024,165DOI:10.1016/J.INFSOF.2023.107336.
\item
  Abufarha S ,Marouf A A ,Rokne G J , et al.Mitigating Prompt Dependency
  in Large Language Models: A Retrieval-Augmented Framework for
  Intelligent Code
  Assistance{[}J{]}.Software,2026,5(1):4-4.DOI:10.3390/SOFTWARE5010004.
\item
  Chang S ,Geng C ,Huang H , et al.CodeSpeak: Improving smart contract
  vulnerability detection via LLM-assisted code analysis{[}J{]}.The
  Journal of Systems \&
  Software,2026,231112635-112635.DOI:10.1016/J.JSS.2025.112635.
\item
  刘昌,周长鑫.BIM模型编码插件的开发与应用{[}J{]}.工程技术研究,2023,8(19):124-126.DOI:10.19537/j.cnki.2096-2789.2023.19.041.
\item
  茆凯成,康彦.基于深度学习技术的智能对话插件开发研究------以心理健康领域为例{[}J{]}.中国新通信,2024,26(24):127-129.
\end{enumerate}

\textbf{致谢}

时光荏苒，本科阶段的学习生涯即将结束。回首这段充实而难忘的求学之路，我在知识探索、实践能力与学术思维上均收获颇丰。在此，我谨向所有给予我指导、支持与陪伴的师长、同学、亲友以及相关单位致以最诚挚的谢意。

首先，我衷心感谢我的指导教师。本论文从选题、框架设计、技术方案到撰写修改，每一步都离不开老师的悉心指引与耐心点拨。老师严谨的治学态度、深厚的学术素养与精益求精的科研精神，使我受益匪浅，不仅顺利完成了毕业设计，更在工程实践与问题解决能力上得到显著提升。

感谢学校为我们提供了良好的学习环境、实验平台与丰富的学术资源，让我能够在专业领域不断探索与实践。感谢所有授课老师四年来的悉心教导，为我打下了坚实的理论基础，使我能够顺利开展研究与开发工作。

感谢同窗好友们在学习与生活中的陪伴与帮助。在项目开发与论文写作过程中，我们相互交流、彼此鼓励、共同进步，这段珍贵的同窗情谊将成为我人生中宝贵的回忆。

最后，我由衷感谢我的家人，他们是我最坚实的后盾，始终给予我无私的关爱、理解与支持，让我能够心无旁骛地完成学业与论文研究。
